
\begin{table}[htbp]
    \centering
    \caption{Identity, Unity, and Individuation: Ontological Distinctions}
    \label{tab:identity-unity-individuation}

    \renewcommand{\arraystretch}{1.4}
    \begin{tabular}{p{3.1cm} p{5.2cm} p{5.3cm}}
        \hline
        \textbf{Conceptual Dimension} &
        \textbf{Ontological Characterization} &
        \textbf{Relation to Identity} \\
        \hline

        Identity &
        The constitutive determination through which an entity is the entity that it is. Identity belongs to the ontological actuality of the entity and is not an additional property. &
        Identity concerns the entity's being itself. It cannot be reduced to unity, numerical singularity, or the conditions through which the entity is distinguished from others. \\

        Unity &
        The ontological coherence through which an entity constitutes a unified whole rather than an arbitrary collection of unrelated elements. &
        Unity concerns the coherence and integration of an entity. It does not, by itself, establish what makes that entity the particular entity it is. \\

        Singularity &
        The condition of being one within a specified determination, domain, or counting relation. &
        Singularity concerns numerical oneness. It must not be treated as a sufficient definition of identity or as proof that an entity possesses only one identity. \\

        Individuation &
        The ontological determination through which an entity is distinguishable as this entity rather than another. &
        Individuation concerns the constitution or establishment of distinction. It is related to identity but does not exhaust the constitutive character of identity. \\

        Numerical identity &
        The relation through which an entity is numerically the same entity across different descriptions, conditions, or temporal positions. &
        Numerical identity is one dimension of the identity problem. It does not independently explain the full ontological constitution of an entity. \\

        Unity of composition &
        The coherence of multiple constituents within an organized entity or whole. &
        Compositional unity may contribute to an entity's constitution, but identity cannot be reduced to the mere aggregation or arrangement of constituents. \\

        Unity of source &
        The unity associated with the ontological source of an entity or a plurality of its manifestations or expressions. &
        Unity of source may be relevant to questions of identity and plurality, but its precise ontological implications require independent examination. \\

        Plurality &
        The existence of multiple entities, determinations, constituents, or manifestations within a specified ontological context. &
        Plurality does not necessarily negate unity or identity. Its relation to identity depends on what is pluralized and at which ontological level. \\

        Multiplicity within an entity &
        The presence of multiple constituents, determinations, relations, or processes within one entity. &
        Internal multiplicity is not sufficient to establish the existence of multiple identities. The relation between multiplicity and identity remains an independent ontological question. \\

        Distinction between entities &
        The ontological differentiation of entities from one another. &
        Distinction is necessary for identifying plurality, but the existence of a distinction does not provide a complete account of the identity of each entity. \\

        Identity and individuation &
        Identity and individuation are related aspects of an entity's ontological constitution, but they address different questions: what the entity is and how it is distinguished as this entity. &
        Neither concept should be defined solely through the other. Their relationship must be established without collapsing constitutive identity into differentiation. \\

        Identity and unity &
        An entity may possess unity through the coherence of its constitution while its identity concerns the determinate character of that entity. &
        Unity and identity are inseparable in certain aspects of an entity's constitution but remain conceptually and analytically distinguishable. \\

        \hline
    \end{tabular}
\end{table}
